\conclusions

Tras el desarrollo integral de esta investigación, se arriba a las siguientes conclusiones fundamentales:

\begin{itemize}
    \item La solución demuestra que es posible captar la variabilidad curricular mediante la 
    técnica de los \ac{FM} de la \ac{SPLE}. El servicio \textit{backend} implementado permite representar formalmente planes de estudio, 
    definir relaciones de obligatoriedad, optatividad y prerrequisitos como restricciones formales, y 
    derivar configuraciones válidas que corresponden a itinerarios formativos personalizados. 
    En consecuencia, se cumple satisfactoriamente el objetivo general planteado al inicio de la investigación.

    \item El servicio \textit{backend} desarrollado se posiciona estratégicamente en el mercado educativo, 
    constituyendo una innovación con pertinencia tecnológica y académico-formativa. Al cubrir la 
    brecha existente entre los sistemas de gestión académica tradicionales 
    (incapaces de modelar variabilidad) y las potentes herramientas de ingeniería de líneas de 
    productos (diseñadas para expertos en software), la solución aporta un enfoque accesible 
    para gestores curriculares no especializados en ingeniería de software, 
    democratizando el uso de técnicas avanzadas de variabilidad en el dominio educativo.
    
    \item La adopción de la metodología ágil \ac{XP} fue acertada al facilitar un desarrollo iterativo y colaborativo, 
    permitiendo una adaptación continua a los requisitos del cliente y una detección temprana de errores. 
    La estrecha participación del usuario durante todo el proceso aseguró que el sistema final se alineara 
    plenamente con las necesidades reales.

    \item La arquitectura en capas adoptada, combinada con los principios de diseño \ac{GRASP} y 
    los patrones \ac{GoF} proporcionó una base sólida para la implementación del sistema. 
    Esta organización permitió una clara separación de responsabilidades, facilitó la mantenibilidad 
    del código y la modularidad del mismo para luego ser sometido a pruebas sin complicaciones.
    
    \item El diagrama de despliegue evidencia que el sistema se despliega sobre un clúster K3s, una 
    distribución ligera de Kubernetes, garantizando alta disponibilidad, escalabilidad automática y tolerancia a fallos. 
    La disposición de los componentes —tres réplicas del \textit{backend} FastAPI, dos trabajadores Celery y 
    los servicios de infraestructura (PostgreSQL, Redis, MinIO) organizados en los espacios de 
    nombres \texttt{ingress} y \texttt{production}— 
    asegura que el tráfico externo entre exclusivamente a través del controlador Traefik, 
    manteniendo el resto de la infraestructura aislada de accesos externos directos. Los 
    volúmenes persistentes y las reclamaciones de volumen persistente garantizan 
    la supervivencia de los datos más allá del ciclo de vida de los contenedores.

    \item Gracias a los resultados obtenidos en las pruebas de rendimiento, el sistema es capaz de mantenerse
    operativo en altas condiciones de demanda, soportando hasta 1000 usuarios simultáneos sin caídas
    del servicio y con tiempos de respuesta inferiores a 3 segundos en el escenario de estrés para las
    operaciones críticas. Los tiempos de respuesta obtenidos garantizan una experiencia de usuario fluida
    incluso durante períodos de planificación curricular intensiva.

\end{itemize} 